\section{Discussion}
\label{sec:discussion}

\subsection{Two regimes and the SDL diagnostic}
SDL separates two regimes of practical consequence. In the
\emph{anchored-contrast} regime, where the target label has an
externally-anchored within-subject neural correlate (EEGMAT
mental-arithmetic alpha desynchronisation; Alzheimer's-disease
resting alpha slowing), fine-tuning projects the label axis onto
the subject-dominated representation and delivers measurable rescue
above classical baselines. In the \emph{contrast-bounded} regime
(UCSD Stress per-recording DASS or longitudinal DSS), no
architecture within five orders of magnitude of parameter scale
exceeds a shared $0.52$--$0.58$ ceiling, and fine-tuning does not
reliably distinguish itself from label-shuffle null. SDL locates the
mechanism behind the aggregate observation from concurrent
benchmarks~\citep{eegbench2025,kauffmann2025critique,eegfmbench2025,
adabrain2025} that EEG FMs sometimes underperform simpler models:
the governing variable is a testable property of the downstream
label (neural-anchor strength), not a training-recipe or scale
artefact. The projection-not-rewrite dynamic we observe also
parallels Chen et al.\ \citep{chen2025sft}'s 757-model LLM
supervised-fine-tuning finding that ``inductive biases of the base
model outweigh the specific SFT corpus''.

\subsection{Pre-benchmark contrast-anchoring protocol}
The operational consequence of SDL is a pre-registration-style
diagnostic that any research group deploying EEG FMs on a new
clinical cohort can apply before investing in benchmark work. Three
steps:
\begin{enumerate}
  \item For the target label, conduct a systematic search for
    published within-subject EEG correlates at the temporal
    resolution of the intended classifier (single-recording,
    trial-level, or trajectory).
  \item If such a literature exists, extract the expected effect size
    and spectral locus and pre-register them as the contrast-strength
    anchor.
  \item If no such literature exists, declare the dataset a
    \emph{contrast-bounded benchmark candidate}, pre-commit to
    subject-level CV plus class-balanced baselines plus
    permutation-null comparison, and report the resulting FM numbers
    as bounded by the subject-dominance ceiling rather than as an
    absolute FM capability claim.
\end{enumerate}
This protocol recasts FM-on-small-clinical-cohort work as a two-stage
enterprise: first anchor the task, then benchmark the architecture.
It inverts the common current practice in which benchmarks are run
first and the absence of a within-subject neural anchor is
discovered, if at all, only when headline numbers fail to replicate
across groups.

\subsection{The TDBRAIN intermediate-anchor case}
\label{sec:discussion_intermediate}
Our paired comparison in Sec.~\ref{sec:results_paired} contrasts a
strong within-subject anchor (EEGMAT alpha-ERD) with an absent
within-subject anchor (UCSD Stress longitudinal DSS). One additional
public dataset in our variance-decomposition analysis,
TDBRAIN \citep{tdbrain} MDD vs HC ($734$ recordings,
$359$ subjects), occupies an intermediate position on the anchor
spectrum under our three-dimension rubric
(Sec.~\ref{sec:methods_anchor}).

\textbf{Empirical anchor result.} Under the same cluster-permutation
protocol used for UCSD Stress, TDBRAIN produces two
cluster-corrected-significant clusters at $p = 0.012$ and $p = 0.014$,
spanning $12.5$--$45$\,Hz over $7$--$8$ fronto-central channels with
within-cluster Hedges' $g \approx 0.4$ (MDD $>$ HC). Whole-scalp
band-averaged effects are beta $g = +0.34$ and gamma $g = +0.33$, with
alpha $g = -0.16$ (weak). This is a different spectral locus from the
published MDD alpha-asymmetry literature the classification was
originally motivated by (broadly traceable to Davidson-tradition
frontal alpha asymmetry); it is directionally consistent with beta/
gamma-band hyperarousal effects reported at cross-subject group level
\citep{olbrich2013mddarousal} but not with a per-recording
within-subject correlate. Full topomap and cluster table are in
App.~A.

\textbf{Interpretation under the three-dimension rubric.} TDBRAIN
satisfies Dimension 2 (empirical cluster survives correction) but
fails Dimension 1 (no published within-subject per-recording
correlate) and Dimension 3 (cluster locus is beta/gamma, not the
published alpha-asymmetry locus). Under our rubric this places
TDBRAIN in the \emph{Intermediate} regime. The three FMs at their
respective best fine-tuning hyperparameters on TDBRAIN MDD reach BAs
of $0.69$ (LaBraM), $0.50$ (CBraMod), and $0.48$ (REVE) under
subject-level cross-validation --- a $19$\,pp spread, substantially
wider than the $5$\,pp band observed on UCSD Stress
(Sec.~\ref{sec:results_paired}).

We read this not as a counterexample to SDL but as an empirical
refinement of the binary anchored/no-anchor classification. In the
strictly-anchored regime the architecture choice is dominated by the
contrast (EEGMAT, ADFTD); in the strictly-bounded regime architecture
is dominated by the subject-dominance ceiling (UCSD Stress); in the
\emph{intermediate-anchor regime} --- where an empirical cluster
exists but misaligns with the published within-subject locus ---
architecture-specific inductive bias appears to dominate in a way
that is neither contrast-driven nor subject-dominance-driven,
producing the observed wide spread. This positions TDBRAIN as a
candidate dataset for the graded-anchor-strength refinement and
points future SDL extensions toward explaining \emph{within-regime}
architectural variance separately from \emph{across-regime} contrast
effects. The Fig.~\ref{fig:sdl_spectrum} spectrum visualises this
position.

\subsection{What SDL does not claim}
We are explicit about three statements SDL does not make. (i) SDL
does not claim that EEG FMs are without value. The anchored-contrast
regime is where EEG FMs deliver meaningful rescue, and this regime
includes well-established clinical targets (dementia, epilepsy,
certain cognitive-load paradigms) whose neural anchors are mature.
(ii) SDL does not claim that the contrast-bounded ceiling is a
permanent limit. A future finding that establishes a robust
per-recording within-subject EEG correlate of DASS or DSS would
change the contrast-bounded classification of UCSD Stress to an
anchored-contrast one, and SDL would correctly predict that FM
fine-tuning would begin to rescue on the newly-anchored task. SDL is
a statement about the current state of the anchoring literature, not
a metaphysical limit on EEG FMs. (iii) SDL does not claim the
two-regime picture is exhaustive. Our controlled paired comparison
spans one strong anchor (alpha-ERD) and one absent anchor (DSS);
intermediate-anchor regimes are outside our tested range and a
productive next step is refining the binary anchor/no-anchor
criterion into a graded anchor-strength score.

\subsection{Implications for pretraining design}
A secondary implication is that SDL suggests pretraining objectives
specifically designed to \emph{reduce} subject-dominance in the
frozen representation --- subject-adversarial pretraining,
subject-invariant contrastive losses, cross-subject covariance
alignment --- would, if successful, shift the contrast-bounded
regime upward even in the absence of an external label anchor.
\citep{multidataset2025} is an early entry in this direction. SDL
predicts that such objectives should show the largest gains
specifically on contrast-bounded benchmarks (where current FMs are
stuck) and smaller gains on contrast-anchored benchmarks (where
current FMs already reach rescue) --- a falsifiable prediction that
future pretraining work can test against our SDL diagnostic.
